\documentclass[11pt,a4paper]{article}
\usepackage[margin=25mm]{geometry}
\usepackage{amsmath,amssymb,amsthm,hyperref}
\setlength{\emergencystretch}{2em}
\newtheorem{lemma}{Lemma}
\newtheorem{theorem}[lemma]{Theorem}
\newcommand{\gen}[1]{\langle#1\rangle}
\title{Routing a maximal shared crossing}
\author{Human Compression project: L044 continuation}
\date{}
\begin{document}\maketitle
Throughout, $G=\gen{\rho_0,\ldots,\rho_{r-1}}\leq A_n$, $n\geq12$,
is transitive, has a connected string C-marking, and every maximal
omission $G_k$ is intransitive. The alternating rank bound is assumed
at strictly smaller degrees at least twelve. All previously accepted
bounded components, including the complete both-imprimitive bound
and the finite small-tail inputs, remain paid.
An eligible $k$-edge crosses two $G_k$-orbits. A unique eligible
edge is a bridge of the full permutation graph and joins its two
omission orbits.

\begin{lemma}[Interval choice in the contradiction range]\label{choice}
Suppose $r>\lfloor(n-1)/2\rfloor$, and a unique eligible $k$-edge
has two nontrivial omission orbits. At least one of their active
index sets is a single interval. If a label larger than $k$ is
shared by both orbits, that interval lies to the right of $k$;
if every shared label is smaller than $k$, it lies to the left.
\end{lemma}
\begin{proof}
If both orbit actions are imprimitive, the complete accepted
both-imprimitive theorem gives the excluded rank bound. Hence at
least one action is primitive. The accepted primitive-orbit interval
lemma gives a single active interval in degree at least three.
Its complete two-point exception theorem supplies the same conclusion
in degree two under the present high-rank hypothesis. This uses
faithful commuting factors and the complete L047 theorem, not IP in
an orbit projection. An active interval avoids the omitted index $k$.
A shared label belongs to both active sets, so its side determines
the side of this interval. In the last assertion a shared label is
understood to exist, as in its applications below.
\end{proof}

\begin{theorem}[A complete numerical maximal-crossing alternative]
\label{routing}
Suppose some unique eligible crossing has a shared label to its
right. Choose its index $i$ maximal among all such crossings.
Then at least one of the following holds:
\begin{enumerate}
\item $r\leq\lfloor(n-1)/2\rfloor$;
\item an actual first or last three-generator subgroup is one of
  the two six-point $A_5$ markings of source Table~2, allowing full
  string reversal;
\item $G_{r-1}$ has a singleton orbit.
\end{enumerate}
The alternatives refer to actual supported subgroups and actual
omission orbits, not projected IP or unnamed isomorphisms.
\end{theorem}
\begin{proof}
Assume the rank bound fails and that no exceptional actual end
triple occurs. We prove the singleton conclusion.
Let the unique $i$-edge be $ab$, joining omission orbits $A,B$,
and let $h>i$ be the last shared label. Both orbits are nontrivial.
For $i>0$, Lemma~\ref{choice} lets us choose $B$ with its active
interval to the right, so $G_{<i}$ fixes $B$ pointwise.
The accepted transitive-side and completed intransitive-side
support alternatives, according to the action of $G_{<i}$ on $A$,
give
\[
 h<r-1,\quad X=\operatorname{supp}(G_{>h})\subseteq B\setminus\{b\},
 \quad 2h\leq n-|X|-1.                                \tag{1}
\]
For $i=0$, the accepted endpoint shared-label exclusion gives
$h<r-1$. All later labels act on just one of $A,B$; they must act
on the same side, since a change of side makes adjacent generators
commute, disconnecting the string diagram. Name that side $B$.
The complete endpoint fixed-support bound gives exactly (1).
Thus (1) holds in both cases.

Put $R=G_{>h}$. If it is transitive on $X$, the complete accepted
transitive-tail alternative contradicts our two exclusions.
Consequently $R$ is intransitive on its support. If every colour
larger than $h$ had two ambient-eligible edges, the accepted ambient
terminal theorem would give
\[
 2(r-1-h)\leq |X|-q\leq|X|-2,
\]
where $q\geq2$ counts its nontrivial orbits. Together with (1)
this gives $2r\leq n-1$, again a contradiction. Every colour has
at least one eligible edge, by transitivity of $G$ and intransitivity
of its omission. Hence there is a smallest $j>h$ with a unique
eligible edge, say $cd$.

Since $\rho_j$ is supported on $X\subset B$, its bridge is inside
$B$. Delete the two bridges. The resulting connected regions are
$A,M,D$, with $a\in A$, $b,c\in M$, $d\in D$; the $G_j$-orbit
containing the $i$-edge is $C=A\cup M$. In particular $C$ is
nontrivial. If $D$ is a singleton and $j<r-1$, the complete interior
singleton-orbit bound contradicts rank failure. If it is a singleton
and $j=r-1$, we have the desired conclusion. We may therefore assume
$D$ nontrivial.

Suppose no label is shared across the $j$-cut. Each generator other
than $\rho_j$ then acts on exactly one of $C,D$. Along either string
interval $<j$ or $>j$ the side cannot change, because adjacent
noncommuting generators cannot have disjoint supports. The prefix
contains the $i$-edge and hence acts on $C$. Nontriviality of $D$
forces $j<r-1$ and makes the suffix act on $D$. Thus these are actual
supported prefix and suffix groups, transitive on their respective
orbits. The complete supported-interval alternative contradicts our
two exclusions. Therefore a shared label exists; let $g$ be the first.
By maximality of $i$, no label larger than $j$ is shared at this
crossing. Thus $g<j$.

We record that $\rho_0$ acts nontrivially on $C$. For $i=0$ its
crossing edge lies in $C$. For $i>0$, it is a nonidentity generator
of $G_{<i}$, supported on $A\subset C$ by our orientation.
This observation fixes the second cut's orientation as follows.

If $j=r-1$, the dual endpoint shared-label exclusion rules out
$g=0$. All indices less than $g$ are unshared and form one connected
string interval. They consequently act on the same side, which is
$C$ because of $\rho_0$. The dual endpoint support bound then gives
\[
 Y=\operatorname{supp}(G_{<g})\subseteq C,\qquad
 2(r-1-g)\leq n-|Y|-1.                                \tag{2}
\]
The suffix $G_{>j}$ is empty in this case.

Suppose instead $j<r-1$. By Lemma~\ref{choice}, one of $C,D$ has
an active interval to the left of $j$. If it were $D$, then
$G_{>j}$ would be supported on $C$. Apply the complete interior
support alternative in reversed index order, with first orbit $C$
and second orbit $D$. Under rank failure it gives $g>0$ and makes
$G_{<g}$ supported on $D$. This is impossible because $\rho_0$
acts on $C$. Therefore the interval side is $C$, and $G_{>j}$ fixes
$C$ pointwise. Apply the same reversed support alternative with
first orbit $D$ and second orbit $C$. It gives $g>0$ and (2).
Both transitive and intransitive actions of this reversed prefix
are covered by the accepted alternatives. We have proved the second
orientation and budget, not assumed them.

Now the complete nested-bridge count applies: its outside regions
$A,D$ are nontrivial, $G_{<i}$ fixes $M\cup D$, and $G_{>j}$ fixes
$A\cup M$. It gives $g\geq h>i$. Since $g$ is the first shared
label at $j$ and $\rho_i$ acts on $C$, this also proves
$\rho_i|_D=1$. Similarly $j>h$ and the last-shared definition at $i$
prove $\rho_j|_A=1$. These are the purity hypotheses needed for the
complete interior and endpoint touching exclusions. Those exclusions
therefore give $g>h$.

Set $P=G_{<g}$. If it is transitive on its support $Y$, the full
transitive-tail alternative in the second cut's reversed orientation
again contradicts the excluded exceptional ends. Thus $P,R$ are
both intransitive on their supports. All colours strictly between
$h$ and $g$ have two ambient-eligible edges, because $g<j$ and $j$
was the first unique crossing after $h$. Equations (1)--(2) and the
complete orbit-incidence deficit theorem now give $2r\leq n-1$.
This final contradiction proves the theorem.
\end{proof}

\paragraph{Source, attribution and exact limit.}
This is a complete replacement of the numerical routing in
Cameron--Fernandes--Leemans--Mixer, \emph{Highest rank of a polytope
for $A_n$}, arXiv:1605.09173v1, Proposition~5.28 (physical pages36--37).
The source's maximal shared crossing and minimal subsequent crossing
strategy is retained. The proper index reversal, degree-two interval
exception, actual orientation, touching exclusions and orbit-incidence
count are supplied by the paid preceding proofs. No new census or
formal declaration is needed.

The literal source conclusion specifies the \emph{last} exceptional
triple; the theorem here retains either actual end, which is the
alternative supplied by the complete transitive-tail theorem.
It does not assert that this end-choice strengthening has been proved.
The main alternating-group recovery still needs the exceptional-head
and singleton continuation, full paired-square rigidity and the exact
remaining source table/disposition assembly. No whole L044, original,
question or priority credit is claimed.
\end{document}
